%%%%%%%%%%%%%%%%%%%%%%%%%%%%%%%%%%%%%%%%%%%%%%%%%%%%%%%%%%%%%%%%%%%%%%
% Lớp 4 | Chủ đề 1. Ôn tập và bổ sung | Ôn tập các phép tính trong phạm vi 100 000
% Dạng 4D1?2-1: Tính nhanh bằng cách nhóm số tròn chục, tròn trăm
%%%%%%%%%%%%%%%%%%%%%%%%%%%%%%%%%%%%%%%%%%%%%%%%%%%%%%%%%%%%%%%%%%%%%%


% nguồn: 2026-11-18_lop4_chu-de-1-tu-luan (nhập 18/11/2026)
\begin{bt}%[4D1N2-1]%[ma:4a407fe7]
% nguon: Toán tư duy lớp 4 – Bài 2 SGK Toán 4 KNTT
	Tính bằng cách thuận tiện: $2\,345 + 4\,500 + 655 + 5\,500$.
	\loigiai{
		Nhóm các số cộng lại được số tròn:\\
		$2\,345 + 4\,500 + 655 + 5\,500 = (2\,345 + 655) + (4\,500 + 5\,500)$\\
		$= 3\,000 + 10\,000 = 13\,000$.
	}
\end{bt}

% nguồn: 2026-11-18_lop4_chu-de-1-tu-luan (nhập 18/11/2026)
\begin{bt}%[4D1H2-1]%[ma:2d8f3a1e]
% nguon: Toán tư duy lớp 4 – Bài 2 SGK Toán 4 KNTT
	Tính bằng cách thuận tiện:\\
	a) $25 \times 7 \times 4$;\qquad b) $5 \times 36 \times 2$.
	\loigiai{
		Nhớ: $25 \times 4 = 100$ và $5 \times 2 = 10$. Ta đổi chỗ các thừa số để nhân hai số này trước.\\
		a) $25 \times 7 \times 4 = (25 \times 4) \times 7 = 100 \times 7 = 700$.\\
		b) $5 \times 36 \times 2 = (5 \times 2) \times 36 = 10 \times 36 = 360$.
	}
\end{bt}

% nguồn: 2026-11-18_lop4_chu-de-1-tu-luan (nhập 18/11/2026)
\begin{bt}%[4D1V2-1]%[ma:caa1962e]
% nguon: Toán tư duy lớp 4 – Bài 2 SGK Toán 4 KNTT
	Tính nhanh: $1 + 2 + 3 + \dots + 19 + 20$.
	\loigiai{
		Ghép số đầu với số cuối thành từng cặp:\\
		$(1 + 20) + (2 + 19) + (3 + 18) + \dots + (10 + 11)$.\\
		Mỗi cặp có tổng bằng $21$. Có $20$ số nên có $20 : 2 = 10$ cặp.\\
		Tổng: $21 \times 10 = 210$.
	}
\end{bt}

% nguồn: 2026-11-18_lop4_chu-de-1-tu-luan-bo-sung (nhập 18/11/2026)
\begin{bt}%[4D1H2-1]%[ma:fbcc9b32]
% nguon: Toán tư duy lớp 4 – Bài 2 SGK Toán 4 KNTT
	Tính bằng cách thuận tiện: $4 \times 125 \times 25 \times 8$.
	\loigiai{
		Ghép các cặp nhân ra số tròn: $4 \times 25 = 100$ và $125 \times 8 = 1\,000$.\\
		$4 \times 125 \times 25 \times 8 = (4 \times 25) \times (125 \times 8) = 100 \times 1\,000 = 100\,000$.
	}
\end{bt}
